\bookchapter{บทที่ 61: เริ่มต้น Swift อย่างรวดเร็ว — ภาษาหลักของการพัฒนาบนแพลตฟอร์ม Apple}

\chapterintro{บทนี้เราจะเรียนภาษากันจริง ๆ แล้ว ไม่ต้องเครียด — Swift มีชื่อเสียงเรื่อง “เป็นมิตรกับมือใหม่” แถมคุณไม่ต้องท่องหนังสือไวยากรณ์: ยึดไวยากรณ์แค่ 20\% ก็อ่านเข้าใจและเขียนโค้ดประจำวันได้ถึง 80\% ที่เหลือค่อยค้นไปทำโปรเจกต์ไป จำหลักการเรียนของเราไว้: เข้าใจ > จำได้ ใช้เป็น > ชำนาญ}

\booksection{ตัวแปรกับชนิดข้อมูล: ตั้งชื่อให้ข้อมูล}

โปรแกรมก็คือ “ข้อมูล + ลอจิก” ตัวแปร ก็คือการตั้งชื่อให้ข้อมูลแล้วเก็บไว้ Swift ประกาศด้วย let กับ var:

\begin{lstlisting}[language=Swift]
let name = "ซานผวน"        // let: ค่าคงที่ (แนะนำให้ใช้เป็นค่าเริ่มต้น)
var age = 18             // var: ตัวแปร (แก้ภายหลังได้)

age = 19                 // <i class="fas fa-check-circle mr-1"></i> var แก้ได้
// name = "ลี่ซือ"         // <i class="fas fa-times-circle mr-1"></i> let แก้ไม่ได้ คอมไพล์แล้ว error ทันที

let score: Int = 100     // จะเขียนไทป์ชัด ๆ ก็ได้
let pi = 3.14            // ไม่เขียนก็เดาเองได้

print("สวัสดี \(name)")   // สอดข้อความ: เครื่องหมายแบ็กสแลช + วงเล็บ
\end{lstlisting}

จำชนิดข้อมูลที่ใช้บ่อยด้วยตารางเดียว:

\begin{center}
\begin{tabularx}{\linewidth}{|X|X|X|}
\hline
\textbf{ชนิดข้อมูล} & \textbf{เก็บอะไร} & \textbf{ตัวอย่าง} \tabularnewline
\hline
Int & จำนวนเต็ม & 42 \tabularnewline
Double & จำนวนทศนิยม & 3.14 \tabularnewline
String & ข้อความ & "สวัสดี" \tabularnewline
Bool & จริง/เท็จ & true / false \tabularnewline
{}[Int] (อาร์เรย์) & ข้อมูลเป็นชุด & {}[1, 2, 3] \tabularnewline
\hline
\end{tabularx}
\end{center}

\begin{center}\small\color{gray}ชนิดข้อมูลพื้นฐานที่ใช้บ่อยใน Swift\end{center}

\begin{tcbinfo}{เคล็ดลับ: หลัก let มาก่อน}
กฎเหล็กของชุมชน Swift: ค่าเริ่มต้นใช้ let เมื่อจำเป็นจริง ๆ ค่อยใช้ var ถ้าจุดที่ใช้ let ได้แต่ใช้ var ไป Xcode จะเตือนด้วยตัวหนังสือสีเหลือง นิสัยนี้ช่วยตัดบั๊กประเภท “ตัวแปรถูกแก้ไปโดยไม่ตั้งใจ” ได้ทั้งตระกูล — โค้ดยิ่งเปลี่ยนน้อย ยิ่งใช้เหตุผลวิเคราะห์ง่าย

\end{tcbinfo}

\booksection{Optional: “เครื่องรางกันแอปพัง” ของ Swift}

นี่คือแนวคิดที่สำคัญที่สุดของ Swift ไม่มีข้อยกเว้น พูดในประโยคเดียว: เมื่อค่าหนึ่ง “อาจไม่มีอยู่” Swift จะบังคับให้คุณจัดการกับมัน — กำจัด “crash จาก null pointer” อาถรรพ์ของวงการนี้ตั้งแต่ระดับภาษา

\begin{lstlisting}[language=Swift]
var name: String = "ซานผวน"   // ไทป์ปกติ: ห้ามเป็น nil เด็ดขาด
// name = nil                // <i class="fas fa-times-circle mr-1"></i> แก้ก่อน ไม่งั้นคอมไพเลอร์ตีกลับทันที

var nickname: String? = nil  // ใส่เครื่องหมายคำถาม: อนุญาตให้เป็น nil (optional)
nickname = "ซานเอ๋อร์"

// ใช้ optional ต้อง "แกะกล่อง" ก่อน:
if let n = nickname {        // if let: แกะแบบปลอดภัย
    print("สวัสดี \(n)")
} else {
    print("ยังไม่มีชื่อเล่น")
}

let len = nickname?.count ?? 0   // ?? ให้ค่าเริ่มต้น (nil coalescing)
print(len)
\end{lstlisting}

\begin{tcbwarn}{ข้อควรรู้: ! คือ “ปฏิบัติการเสี่ยงอันตราย”}
เขียน nickname! ในโค้ด (force unwrap) เท่ากับตะโกนบอกคอมไพเลอร์ว่า “ฉันรับประกันว่ามันไม่ใช่ nil” — พอเดาผิด แอปพังต่อหน้าต่อตา หลักช่วงมือใหม่: พยายามใช้ if let / ?? อย่าใช้ ! บางที Xcode เติมโค้ดให้แล้วแถม ! มาด้วย อย่าลืมเอาออก

\end{tcbwarn}

\booksection{ฟังก์ชันกับ closure: ห่อลอจิกเป็น “ปุ่มกด”}

ฟังก์ชันคือลอจิกที่มีชื่อ อย่างปุ่มบนรีโมต ส่วน closure คือ “ฟังก์ชันไร้ชื่อ” ที่หยิบยื่นให้คนอื่นได้ทันที:

\begin{lstlisting}[language=Swift]
// ฟังก์ชัน: ชื่อมาก่อน พารามิเตอร์และไทป์คืนค่ามาหลัง
func greet(name: String) -> String {
    return "สวัสดี \(name)!"
}
print(greet(name: "ซานผวน"))     // ตอนเรียกต้องระบุชื่อพารามิเตอร์

// คลูเซอร์: นิยามเลยตรงนั้น กระชับไว
let numbers = [1, 2, 3, 4]
let doubled = numbers.map { $0 * 2 }   // $0 = ตัวปัจจุบัน
// [2, 4, 6, 8]

let big = numbers.filter { $0 > 2 }    // กรอง
// [3, 4]
\end{lstlisting}

ใน SwiftUI closure อยู่ทุกหนแห่ง: วงเล็บปีกกาหลัง Button(“กดฉัน”) \{ ... \} ก็คือ closure หนึ่ง map, filter, sorted ท่าคอมโบพวกนี้จำให้แม่น — ต่อจากนี้การประมวลผลลิสต์ 90\% ล้วนเป็นการผสมกันของพวกมัน

\booksection{struct กับ protocol: ระบบ “เลโก้” ของ Swift}

จุดยืนทางการของ Swift ชัดเจน: ให้ใช้ struct เป็นหลัก ค่อยใช้ class เมื่อต้องการ inheritance โดย struct เป็น “value type” — พอคัดลอกก็แยกกันไปใช้ชีวิตคนละทาง ตัดบั๊กได้ทั้งประเภทโดยธรรมชาติ:

\begin{lstlisting}[language=Swift]
// struct: กล่องใส่ข้อมูล
struct Record {
    var id = UUID()
    var amount: Double
    var category: String
    var date: Date

    // struct มี "memberwise initializer" ให้พร้อม: สร้างได้เลย
}
let r = Record(amount: 25.5, category: "อาหาร", date: Date())

// protocol: สัญญาว่า "ต้องมีความสามารถอะไรบ้าง"
protocol Identifiable {
    var id: UUID { get }
}
struct Record: Identifiable { ... }   // เชื่อฟัง protocol = สัญญาว่าจะทำตามที่กำหนด
\end{lstlisting}

ใน SwiftUI แทบทุกอย่างเป็น struct (View ก็คือ protocol) ประโยคนี้ตอนนี้ไม่เข้าใจก็ไม่เป็นไร บทที่ 62 จะได้เจอกันทุกวัน — ในโลกของ Swift ให้กอด struct ไว้ก่อน

\booksection{การ์ดค้นเร็วไวยากรณ์สำหรับคนเริ่มจากศูนย์}

\begin{center}
\begin{tabularx}{\linewidth}{|X|X|}
\hline
\textbf{ต้องการทำอะไร} & \textbf{วิธีเขียน} \tabularnewline
\hline
ประกาศค่าคงที่/ตัวแปร & \texttt{let x = 1}/\texttt{var x = 1} \tabularnewline
แทรกค่าลงในสตริง & "สวัสดี \textbackslash{}(name)" \tabularnewline
ประกาศฟังก์ชัน & func ชื่อฟังก์ชัน(พารามิเตอร์: ชนิดข้อมูล) -> ชนิดค่าที่คืน \{ ... \} \tabularnewline
การเช็คเงื่อนไข & if เงื่อนไข \{ ... \} else \{ ... \} \tabularnewline
การวนลูป & for x in อาร์เรย์ \{ ... \} \tabularnewline
การคลายค่าอย่างปลอดภัย & if let x = ค่าทางเลือก \{ ... \} \tabularnewline
ค่าเริ่มต้น & x ?? ค่าเริ่มต้น \tabularnewline
การนิยาม struct & \texttt{struct X \{ var a: Int; var b: String \}} \tabularnewline
Closure & \texttt{\{ \$0 * 2 \}} \tabularnewline
\hline
\end{tabularx}
\end{center}

\begin{center}\small\color{gray}บัตรค้นเร็วไวยากรณ์ Swift ฉบับมือใหม่\end{center}

\begin{tcbpractice}{ลงมือทำ: เล่น Swift ใน Playgrounds 15 นาที}
มี Mac ก็เปิด Xcode Playground; ไม่มี Mac ก็ใช้ Swift Playgrounds บน iPad หรือสนามซ้อมเวอร์ชันเว็บ รันโค้ดทุกช่วงในบทนี้ให้ครบ ลองแก้ตัวเลขและข้อความดู และตั้งใจทำให้เกิด error สักสองสามครั้งเพื่อดูว่า Xcode จะแจ้งเตือนอย่างไร 15 นาทีของการ'เล่นก่อนเรียน'แบบนี้ จะทำให้คุณเริ่มต้นบทที่ 62 ได้ง่ายขึ้นเท่าตัว

\end{tcbpractice}

\begin{tcbsummary}{สรุปบทนี้}
\begin{itemize}
  \item ใช้ let เป็นหลัก ส่วน var ใช้อย่างระมัดระวัง; ชนิดข้อมูลอนุมานให้เอง หรือจะระบุชัดหลังเครื่องหมายทวิภาค (:) ก็ได้
  \item Optional คือจิตวิญญาณของ Swift: ค่าที่อาจว่างต้องจัดการให้ครบ if let / ?? คือเครื่องรางของขลัง ส่วน ! คือการเสี่ยงระดับสูง
  \item การเรียกฟังก์ชันต้องระบุชื่อพารามิเตอร์; Closure คือฟังก์ชันไร้ชื่อ ที่พบได้ทั่วไปใน SwiftUI
  \item ใช้ struct เป็นหลัก และคอนฟอร์มกับ Protocol; แทบทุกอย่างใน SwiftUI ล้วนเป็น struct
  \item ไวยากรณ์แค่ 20\% ทำงานได้ถึง 80\%: ใช้ให้เป็นก่อน แล้วค่อยเชี่ยวชาญ
\end{itemize}

\end{tcbsummary}
